\documentclass[a4paper,12pt]{article}
\usepackage[utf8]{inputenc}
\usepackage{amsmath}
\usepackage{listings}
\usepackage{xcolor}
\usepackage{hyperref}
\usepackage{geometry}
\geometry{margin=1in}

\title{The mighf ISA: A Next-Gen Minimalist Instruction Set Architecture}
\author{Miguel “Kuyo” / SufremOak}
\date{\today}

% Code highlighting style for assembly
\lstdefinelanguage{mighfasm}{
    keywords={MOV, ADD, SUB, JMP, CALL, RET, PUSH, POP, LOAD, STORE, AND, OR, XOR, NOT, SHL, SHR, NOP, HALT},
    keywordstyle=\color{blue}\bfseries,
    comment=[l]{;},
    commentstyle=\color{gray}\itshape,
    morecomment=[s]{/*}{*/},
    morekeywords=[2]{R0,R1,R2,R3,R4,R5,R6,R7,PC,SP,FLAGS},
    keywordstyle=[2]\color{teal},
    sensitive=true,
    basicstyle=\ttfamily\small,
    numbers=left,
    numberstyle=\tiny\color{gray},
    stepnumber=1,
    numbersep=5pt,
    backgroundcolor=\color{white},
    showspaces=false,
    showstringspaces=false,
    tabsize=4,
    breaklines=true,
    breakatwhitespace=true,
}

\begin{document}

\maketitle

\begin{abstract}
The mighf Instruction Set Architecture (ISA) is a custom-designed, minimal yet powerful ISA aimed at embedded and low-level system programming. It balances simplicity and extensibility while supporting modern needs like stack operations, conditional branching, and a basic type system.
\end{abstract}

\section{Introduction}
The mighf ISA is designed for efficient assembly-level programming on low-resource machines, blending classic RISC simplicity with flexible stack-based operations. It is intended to be a clean foundation for multiple platform ports and language toolchains.

\section{Register Set}
\begin{itemize}
    \item \textbf{R0--R7}: General purpose registers.
    \item \textbf{PC}: Program Counter — holds the address of the next instruction.
    \item \textbf{SP}: Stack Pointer — points to the top of the stack.
    \item \textbf{FLAGS}: Status flags register (Zero, Carry, Sign, Overflow).
\end{itemize}

\section{Memory Model}
The ISA uses a flat 16-bit address space (adjustable for target platforms). Memory is byte-addressable. The stack grows downward in memory.

\section{Instruction Formats}
The ISA supports fixed 16-bit instruction encoding with optional immediate values and register operands.

\section{Core Instructions}

\begin{lstlisting}[language=asm,caption={Example mighf assembly code}]
; Move immediate value 5 into R1
MOV R1, #5        

; Add R1 and R2, store in R3
ADD R3, R1, R2   

; Push R3 onto stack
PUSH R3          

; Conditional jump if zero flag set
JZ label_zero     

; Return from subroutine
RET              

label_zero:
NOP               
HALT              
\end{lstlisting}

Key instructions include:
\begin{itemize}
    \item \texttt{MOV} — Move data between registers or immediate to register.
    \item \texttt{ADD}, \texttt{SUB}, \texttt{AND}, \texttt{OR}, \texttt{XOR} — Basic ALU ops.
    \item \texttt{PUSH}, \texttt{POP} — Stack operations.
    \item \texttt{JMP}, \texttt{JZ}, \texttt{JNZ} — Control flow.
    \item \texttt{CALL}, \texttt{RET} — Subroutine call and return.
    \item \texttt{LOAD}, \texttt{STORE} — Memory access.
\end{itemize}

\section{Flags}
\begin{itemize}
    \item \textbf{Z} (Zero): Set if the result of an operation is zero.
    \item \textbf{C} (Carry): Set if an operation results in a carry out.
    \item \textbf{S} (Sign): Set if the result is negative.
    \item \textbf{O} (Overflow): Set if arithmetic overflow occurs.
\end{itemize}

\section{Extensibility}
The ISA is designed to allow future extensions including new instructions, wider data types, and custom coprocessor operations.

\section{Conclusion}
mighf ISA strikes a balance between simplicity and practical features needed for modern embedded programming. It is a solid base for further tooling and OS development.

\end{document}
